% from ega4-2-fr.tex
\oldpage[IV-2]{165}
de générateurs $t'_j$ $(1\leq j\leq m)$ de $K'$ sur
$K(t_1,\ldots,t_n)$, que l'on peut supposer entiers sur
$A_1=A[t_1,\ldots,t_n]$, on voit que
$A'_1=A_1[t'_1,\ldots,t'_m]$ est fini sur $A_1$ et a pour corps de fractions
$K'$. Si on pose $X'=\operatorname{Spec}(A')$,
$X'_1=\operatorname{Spec}(A'_1)$, il résulte de ce que les corps de fonctions
rationnelles $R(X')$ et $R(X'_1)$ sont tous deux isomorphes à $K'$ et de ce
que $X'$ et $X'_1$ sont des $A$-préschémas de type fini, qu'il existe un
ouvert $U'\subset X'$ et un ouvert $U'_1\subset X'_1$ qui sont
$A$-isomorphes (I, 6.5.5). On est donc ramené à prouver que
$\operatorname{Reg}(X'_1)$ contient un ouvert non vide, autrement dit on
peut supposer que $A'$ est une $A_1$-algèbre finie. Or on sait (0, 17.3.7)
que $A_1$ est un anneau régulier, et l'on peut donc se borner au cas où $A'$
est une $A$-algèbre finie et $K'$ une extension séparable et finie de $K$.
Si $\xi$ est le point générique de $X$, $A'_\xi=K'$ est alors un module
libre sur $A_\xi=K$, donc (Bourbaki, \emph{Alg. comm.}, chap. II, \S5,
n\textsuperscript{o} 1,
cor. de la prop. 2) on peut, en remplaçant au besoin $A$ par un $A_f$,
supposer que $A'$ est un $A$-module libre de type fini. Soit alors
$(x_h)_{1\leq h\leq r}$ une base de ce $A$-module, et posons
\[
d=\det(\operatorname{Tr}_{A'/A}(x_hx_k))=
\det(\operatorname{Tr}_{K'/K}(x_hx_k))\in A.
\]
Comme $K'$ est séparable sur $K$, on sait (Bourbaki, \emph{Alg.}, chap. IX, \S2,
prop. 5) que $d\neq0$; remplaçant au besoin $A$ par l'anneau de fractions
$A_d$, on peut supposer $d$ inversible dans $A$. Mais alors, pour tout
$z\in\operatorname{Spec}(A)$, si l'on désigne par $\bar x_h$
$(1\leq h\leq r)$ l'image canonique de $x_h$ dans
$A'(z)=A'\otimes_A k(z)$, on a
$\det(\operatorname{Tr}_{A'(z)/k(z)}(\bar x_h\bar x_k))=\bar d$, image
canonique de $d$ dans $k(z)=A_z/\mathfrak m_z$ et comme $\bar d$ est
inversible (donc $\neq0$) dans $k(z)$, on sait (loc. cit.) que $A'(z)$ est
une $k(z)$-algèbre séparable, donc composée directe de corps extensions
finies et séparables de $k(z)$. Une telle algèbre étant un anneau régulier,
on voit que le morphisme $g:X'\to X$ est plat et que ses fibres $g^{-1}(z)$
sont régulières pour tout $z\in X$; on conclut alors de (6.5.2, (ii)) que
$X'$ est régulier, ce qui termine dans ce cas la démonstration.
\end{proof}

II) Cas général. --- Comme $A'$ est un $A$-module sans torsion,
$A'\otimes_AK$ s'identifie à un sous-anneau de $K'$, donc
$X''=\operatorname{Spec}(A'\otimes_AK)$ est un $K$-préschéma intègre, dont
$K'$ est le corps des fonctions rationnelles. On sait (4.6.6) qu'il existe
une extension radicielle finie $K_1$ de $K$ telle que si
$X'_1=\operatorname{Spec}(A'\otimes_AK_1)=X''\otimes_KK_1$,
$(X'_1)_{\mathrm{red}}$ soit un $K_1$-préschéma géométriquement réduit de
type fini; en outre, le morphisme $\operatorname{Spec}(K_1)\to
\operatorname{Spec}(K)$ étant radiciel, fini et surjectif, est un
homéomorphisme universel (2.4.5), donc $X'_1$ est homéomorphe à $X''$, et
par suite $(X'_1)_{\mathrm{red}}$ est intègre; en outre, son corps de
fonctions rationnelles $K'_1$ est une extension radicielle finie de $K'$ et
une extension séparable de type fini de $K_1$ (4.6.1). En vertu de
l'hypothèse c) de l'énoncé, il y a une sous-$A$-algèbre finie $A_1$ de
$K_1$, ayant $K_1$ pour corps des fractions et telle que si l'on pose
$X_1=\operatorname{Spec}(A_1)$, $\operatorname{Reg}(X_1)$ contienne un
ouvert non vide de $X_1$. Soit $A'_1$ l'image de l'homomorphisme canonique
$A'\otimes_AA_1\to K'_1$ et soit $X'_1=\operatorname{Spec}(A'_1)$; $A'_1$
est un anneau intègre qui est une $A'$-algèbre finie et dont le corps des
fractions est $K'_1$ par construction; en outre, comme l'homomorphisme
composé $A'\to A'\otimes_AA_1\to A'_1\to K'_1$ est identique à
$A'\to K'\to K'_1$, donc injectif, l'homomorphisme $A'\to A'_1$ est
injectif; le morphisme $g:X'_1\to X'$ est donc fini et surjectif
(II, 6.1.10). Cela étant, l'hypothèse sur $A_1$ et la partie I) de la
démonstration entraînent que $\operatorname{Reg}(X'_1)$
